\begin{animation}[id="persistent-segtree", label="Persistent Segment Tree -- Cap nhat sum that qua value-layer"]
\shape{T}{Tree}{data=[1,2,3,4,5,6,7,8], kind="segtree", show_sum=true, label="Segment tree — sums hien tren moi node"}

\step
\apply{T.node["[0,7]"]}{value="36"}
\apply{T.node["[0,3]"]}{value="10"}
\apply{T.node["[0,1]"]}{value="3"}
\apply{T.node["[0,0]"]}{value="1"}
\apply{T.node["[1,1]"]}{value="2"}
\apply{T.node["[2,3]"]}{value="7"}
\apply{T.node["[2,2]"]}{value="3"}
\apply{T.node["[3,3]"]}{value="4"}
\apply{T.node["[4,7]"]}{value="26"}
\apply{T.node["[4,5]"]}{value="11"}
\apply{T.node["[4,4]"]}{value="5"}
\apply{T.node["[5,5]"]}{value="6"}
\apply{T.node["[6,7]"]}{value="15"}
\apply{T.node["[6,6]"]}{value="7"}
\apply{T.node["[7,7]"]}{value="8"}
\narrate{\textbf{Phien ban V0}: segment tree xay tu mang \texttt{[1,2,3,4,5,6,7,8]}. Moi node hien tong cua doan $[lo, hi]$ ma no dai dien. Tong cua ca mang = 36 (tai node \texttt{[0,7]}). Day la trang thai goc ma persistent segment tree se giu nguyen mai mai, bat ke co bao nhieu ban cap nhat sau nay.}

\step
\recolor{T.node["[3,3]"]}{state=current}
\recolor{T.node["[2,3]"]}{state=current}
\recolor{T.node["[0,3]"]}{state=current}
\recolor{T.node["[0,7]"]}{state=current}
\recolor{T.node["[0,1]"]}{state=dim}
\recolor{T.node["[0,0]"]}{state=dim}
\recolor{T.node["[1,1]"]}{state=dim}
\recolor{T.node["[2,2]"]}{state=dim}
\recolor{T.node["[4,7]"]}{state=dim}
\recolor{T.node["[4,5]"]}{state=dim}
\recolor{T.node["[4,4]"]}{state=dim}
\recolor{T.node["[5,5]"]}{state=dim}
\recolor{T.node["[6,7]"]}{state=dim}
\recolor{T.node["[6,6]"]}{state=dim}
\recolor{T.node["[7,7]"]}{state=dim}
\apply{T.node["[3,3]"]}{value="14"}
\apply{T.node["[2,3]"]}{value="17"}
\apply{T.node["[0,3]"]}{value="20"}
\apply{T.node["[0,7]"]}{value="46"}
\narrate{\textbf{V1 — cap nhat $a[3]$ += 10}. Chi cac node tren \textbf{path-copy} tu la \texttt{[3,3]} len root moi phai tao ban sao moi: \texttt{[3,3]} 4 → 14, \texttt{[2,3]} 7 → 17, \texttt{[0,3]} 10 → 20, \texttt{[0,7]} 36 → 46. Cac node dim van chia se voi V0 (nguyen ven) — persistent khong nhan ban nhung node khong doi. Chi phi: $O(\log n)$ node moi moi cap nhat. Day la cap nhat \textbf{that}: value-layer ghi de len label tinh, nen so hien thi thay doi qua tung frame.}

\step
\recolor{T.node["[3,3]"]}{state=good}
\recolor{T.node["[2,3]"]}{state=good}
\recolor{T.node["[0,3]"]}{state=good}
\recolor{T.node["[0,7]"]}{state=current}
\recolor{T.node["[5,5]"]}{state=current}
\recolor{T.node["[4,5]"]}{state=current}
\recolor{T.node["[4,7]"]}{state=current}
\recolor{T.node["[0,1]"]}{state=dim}
\recolor{T.node["[0,0]"]}{state=dim}
\recolor{T.node["[1,1]"]}{state=dim}
\recolor{T.node["[2,2]"]}{state=dim}
\recolor{T.node["[4,4]"]}{state=dim}
\recolor{T.node["[6,7]"]}{state=dim}
\recolor{T.node["[6,6]"]}{state=dim}
\recolor{T.node["[7,7]"]}{state=dim}
\apply{T.node["[5,5]"]}{value="11"}
\apply{T.node["[4,5]"]}{value="16"}
\apply{T.node["[4,7]"]}{value="31"}
\apply{T.node["[0,7]"]}{value="51"}
\narrate{\textbf{V2 — cap nhat $a[5]$ += 5} tiep tu V1. Path moi la \texttt{[5,5]} → \texttt{[4,5]} → \texttt{[4,7]} → \texttt{[0,7]}. Gia tri moi: 6 → 11, 11 → 16, 26 → 31, 46 → 51. Luu y root \texttt{[0,7]} la node duy nhat nam tren \textbf{ca hai} path nen gia tri cua no thay doi cho lan thu hai (46 → 51). Cac node trong nhanh V1 cu (\texttt{[3,3], [2,3], [0,3]}) gio duoc to mau \texttt{good} — chung thuoc V1 va van duoc giu nguyen ven cho truy van V1.}

\step
\recolor{T.node["[3,3]"]}{state=good}
\recolor{T.node["[2,3]"]}{state=good}
\recolor{T.node["[0,3]"]}{state=good}
\recolor{T.node["[5,5]"]}{state=good}
\recolor{T.node["[4,5]"]}{state=good}
\recolor{T.node["[4,7]"]}{state=good}
\recolor{T.node["[0,7]"]}{state=good}
\recolor{T.node["[0,1]"]}{state=done}
\recolor{T.node["[0,0]"]}{state=done}
\recolor{T.node["[1,1]"]}{state=done}
\recolor{T.node["[2,2]"]}{state=done}
\recolor{T.node["[4,4]"]}{state=done}
\recolor{T.node["[6,7]"]}{state=done}
\recolor{T.node["[6,6]"]}{state=done}
\recolor{T.node["[7,7]"]}{state=done}
\narrate{Tong ket: ba phien ban cung ton tai. V0 co root=36 (mang goc), V1 co root=46 ($a[3]=14$), V2 co root=51 ($a[3]=14, a[5]=11$). Truy van bat ky phien ban nao chi ton $O(\log n)$. Bo nho tong: $O(n + q \log n)$ thay vi $O(nq)$ nho viec chia se cac node khong nam tren path. Ung dung: k-th smallest trong range (Chairman tree), offline queries, version rollback, time-travel data structures.}
\end{animation}
